\scanpage{39}
\margin{$/^{\circ}$ NB. S'il s'agit de \unsure{grpes|groupes} de type \unsure{galoisien|quelconque|galoisiens} et $U$ est ouvert \unsure{en|sur|dans} $I$, alors il s'ensuit qu'il contient un ss-groupe ouvert \unsure{plus|fin|de} $U_{\unsure{1|0}}$, qui est \uline{invariant} sous $D$}
\uline{\unsure{Sorite|Suite|Boîte} de la \unsure{moyenne|remarque}}

$D$ groupe, $I$ \unsure{ss-groupe|sous-groupe} distingué, $E$ vectoriel de dim finie sur corps $k$, sur lequel $D$ opère. On suppose que $I$ admet un sous-groupe d'indice fini $U$, \inserted{invariant \uline{dans} $D/^{\circ}$,} tel que $g\in U\Rightarrow$ $g$ opère sur $E$ de façon \uline{unipotente}, i.e. valeurs propres \unsure{ttes|toutes} $1$. Alors le groupe algébrique d'autom de $E$ engendré par les $g\in U$ est ``unipotent'', donc \deleted{\unsure{si} \unsure{on est|on a}} $E$ admet une \unsure{suite} de composition sur $U$ tel que \deleted{les} \inserted{$E$ \unsure{opère|quot.} trivialement sur les} facteurs. \deleted{\unsure{soient|sont|soit} stables sous $E$}. De façon précise, on peut prendre la suite de composition par les $E^U$, $(E/E^U)^U$ .... . Cette suite est stable par les opérations de $D$, \uline{si} \deleted{(comme il est loisible)} \deleted{on prend} \inserted{puisque} $U$ dans $I$ invariant \unsure{pour|par} les $D$. Donc il existe une suite de composition de $E$ \uline{sur $D$} tel que sur chaque facteur $E_i$, $\deleted{\unsure{I|E}}\,U$
% archive sheet number: 38
